\documentclass[a4paper,12pt]{article}

% --- 宏包引入 ---
\usepackage{ifpdf}
\usepackage{ifxetex}
\usepackage{ifluatex}
\ifxetex
  \usepackage{xeCJK}
  \setCJKmainfont{SimSun} % 设置中文字体，如宋体
\else
  \ifluatex
    \usepackage{fontspec}
  \else
    \usepackage[UTF8]{ctex}
  \fi
\fi

% --- 页面布局 ---
\usepackage{geometry}
\geometry{a4paper, left=2.5cm, right=2.5cm, top=2.5cm, bottom=2.5cm}
\linespread{1.5}
% --- 其它宏包 ---
\usepackage{titlesec}       % 标题格式
\usepackage{array}          % 表格扩展
\usepackage{booktabs}       % 三线表
\usepackage{multirow}       % 表格合并行
\usepackage{graphicx}       % 图片插入
\usepackage{float}          % 图片浮动位置控制
\usepackage{amsmath}        % 数学公式
\usepackage{amssymb}        % 数学符号
\usepackage{xcolor}         % 颜色支持
\usepackage{fancyhdr}       % 页眉页脚
% 代码环境设置
\usepackage{listings}
\lstset{
    basicstyle=\ttfamily\small,
    breaklines=true,
    frame=single,
    numbers=left,
    numberstyle=\tiny,
    showspaces=false,
    showstringspaces=false,
    tabsize=4,
    language=Python
}
\pagestyle{fancy}     % 设置页面风格为 fancy
\fancyhf{}            % 清空默认的页眉页脚
\setlength{\headwidth}{\textwidth}
\lhead{人工智能基础 A 综合实践实验报告} % 设置居中页眉 (c=center, head=页眉)
\rhead{蒋贤迪 3240105782}

% --- 标题设置 ---
\titleformat{\section}{\normalfont\large\bfseries}{\thesection}{1em}{}
\titlespacing*{\section}{0pt}{3.5ex plus 1ex minus .2ex}{2.3ex plus .2ex}
\title{\textbf{\fontsize{18pt}{22pt}\selectfont 人工智能基础A综合实践实验报告}}
\author{}
\date{}

% --- 正文开始 ---
\begin{document}

\vspace{-2em}

% 标题部分
\maketitle
\thispagestyle{fancy} % 首页不显示页码

\vspace{-3cm}

% --- 一、基本信息 (仿照大作业.pdf 表格) ---
\section*{一、基本信息}

\begin{table}[H]
    \centering
    \renewcommand{\arraystretch}{1.6} % 表格行高调整
    \setlength{\tabcolsep}{8pt}       % 列间距调整
    \begin{tabular}{|c|c|l|c|l|}
        \hline
        \multicolumn{2}{|c|}{\textbf{是否组队}} & \multicolumn{1}{l|}{$√$ 是 \quad $\square$ 否} & \multicolumn{1}{c|}{\textbf{组名}} & \multicolumn{1}{l|}{碱基互补配队} \\ 
        \hline
        \multirow{3}{*}{\textbf{成员信息}} & \textbf{成员1} & \multicolumn{3}{l|}{姓名：\underline{蒋贤迪} \quad 学号：\underline{3240105782}} \\ \cline{2-5}
         & \textbf{成员2} & \multicolumn{3}{l|}{姓名：\underline{\hspace{4cm}} \quad 学号：\underline{\hspace{4cm}}} \\ \cline{2-5}
         & \textbf{成员3} & \multicolumn{3}{l|}{姓名：\underline{\hspace{4cm}} \quad 学号：\underline{\hspace{4cm}}} \\
        \hline
        \multicolumn{2}{|c|}{\textbf{选择框架}} & \multicolumn{3}{l|}{$√$ PyTorch \quad $\square$ TensorFlow \quad $\square$ scikit-learn \quad $\square$ 飞桨 \quad $\square$ 其他：\underline{\hspace{1cm}}} \\
        \hline
        \multicolumn{2}{|c|}{\textbf{选题}} & \multicolumn{3}{l|}{基于 RNA-FM 与 RNA-ret 的 RNA 二级结构预测平台 (RNA-Dee)} \\
        \hline
        \multicolumn{2}{|c|}{\textbf{数据集}} & \multicolumn{3}{l|}{$\square$ 中药数据集 \quad $√$ 其他：\underline{\hspace{0.5cm} RNAStrAlign \hspace{0.5cm}}} \\
        \hline
        \multicolumn{2}{|c|}{\textbf{基准测试集}} & \multicolumn{3}{l|}{ArchiveII} \\
        \hline
    \end{tabular}
    \vspace{0.5em}
\end{table}

\noindent \textbf{数据集说明：}

本实验采用RNAStrAlign和ArchiveII数据集。该数据涵盖了 tRNA, 5S rRNA, 16S rRNA 等多个 RNA 家族。
\begin{itemize}
    \item \textbf{数据格式}：原始数据为 \texttt{.bpseq} 格式（包含 Index, Residue, Pair Index 三列），经预处理脚本解析为序列与 2D 接触矩阵。
    \item \textbf{预处理}：设置最大序列长度为 1022，去除过长或无效序列。
    \item \textbf{数据划分}：数据集按 7:3 划分为训练集和验证集，固定种子为42。测试集采用独立的 ArchiveII 数据集进行最终基准评估。
    \item \textbf{测试集说明}：测试集不仅包含了训练集里RNA的种类，还有其他种类的RNA，如23sRNA等，能够有效评估模型的泛化能力。
\end{itemize}

\newpage

% =========================================================
% 二、评估设置 (严格复刻 PDF 原文)
% =========================================================
\section*{二、评估设置}

\subsection*{2.1 评估指标选择}
\noindent 本实验选择了准确率、精确率、召回率 和 F1 分数 作为评估指标。它们的计算公式如下：

\begin{equation}
\begin{aligned}
\text{准确率}Accuracy &= \frac{TP + TN}{TP + TN + FP + FN} \\
\text{精确率}Precision &= \frac{TP}{TP + FP} \\
\text{召回率}Recall &= \frac{TP}{TP + FN} \\
F1 &= 2 \cdot \frac{Precision \cdot Recall}{Precision + Recall}
\end{aligned}
\end{equation}

\noindent \textbf{符号说明：}
\begin{itemize}
    \item \textbf{TP (True Positive)}: 真实结构中存在碱基配对，且模型正确预测为配对。
    \item \textbf{TN (True Negative)}: 真实结构中无配对，且模型正确预测为无配对。
    \item \textbf{FP (False Positive)}: 真实结构中无配对，但模型错误预测为配对。
    \item \textbf{FN (False Negative)}: 真实结构中存在配对，但模型错误预测为无配对。
\end{itemize}

\vspace{1em}
\noindent \textbf{采用这些评估指标的原因：}

\begin{enumerate}
    \item \textbf{应对数据类别不平衡}：
    
    \hspace{2em}RNA 二级结构预测生成的接触图是一个高度稀疏的矩阵。在 $L \times L$ 的矩阵中，绝大多数位置是未配对的（0），只有极少数位置存在碱基配对（1）。在这种情况下，如果模型倾向于预测全 0，虽然能获得极高的\textbf{准确率}，但完全失去了预测结构的意义。因此，单纯依赖准确率不足以客观评估模型性能。

    \item \textbf{综合考量查准与查全}：
    
    \hspace{2em}精确率反映了模型预测出的配对有多少是可信的，而\textbf{召回率} 反映了模型找出了多少真实的配对结构。为了平衡这两者，本实验采用 F1 分数 作为核心评价指标。F1 分数是精确率和召回率的调和平均数，只有当两者都较高时，F1 分数才会高，这最能体现模型在稀疏矩阵预测任务中的综合表现。
\end{enumerate}

\newpage

\subsection*{2.2 测试数据集及划分标准}


\hspace{2em}依据实验要求，测试数据集应完全独立于模型训练与调参过程。本实验采用了 ArchiveII 数据集作为最终的基准测试集。

\begin{itemize}
    \item \textbf{独立性说明}：
    
    \hspace{2em}ArchiveII 数据集中的序列与训练集（RNAStrAlign）虽有部分生物学上的同源性，但在本实验中，ArchiveII 被严格隔离，从未参与训练或验证阶段的任何计算。
    \item \textbf{评估目的}：
    
    \hspace{2em}使用该数据集旨在客观评估模型在未见过的真实 RNA 序列上的预测性能，验证模型的泛化鲁棒性。
\end{itemize}
% =========================================================
% 三、代码实现
% =========================================================
\section*{三、代码实现}
\subsection*{3.1 基于预训练模型的特征提取}

\hspace{2em}本实验的核心创新之一在于摒弃了传统的独热编码（One-hot）方式，转而采用在大规模未标注数据上进行过自监督学习的 \textbf{RNA-FM} 模型作为特征提取器。该模型能够捕捉序列内部深层的进化语义信息与上下文依赖关系。

\noindent \textbf{1. 模型加载与初始化}

我们在代码中通过调用接口加载预训练权重。为了保证特征提取器的稳定性，我们在初始化阶段冻结了骨干网络的全部参数，仅将其作为静态的特征生成器使用。

\noindent \textbf{2. 第12层特征截取}

根据相关文献研究，预训练模型的深层输出通常包含更丰富的结构化信息。因此，本实验并未直接使用模型的最终输出，而是提取了第 12 层变换器（Transformer）的隐藏状态作为下游任务的输入。

\noindent \textbf{3. 维度说明}

对于长度为 $L$ 的输入序列，经过分词处理后进入模型，提取出的特征矩阵维度为 $L \times 640$。其中 640 维的向量代表了每个碱基在进化空间中的高维语义表示。

\begin{figure}[H]
    \centering
    \begin{minipage}{0.85\linewidth}
        \centering
        \includegraphics[width=\textwidth]{代码实现/RNA-FM.png} % 请替换为实际截图文件名
        \caption{特征提取模块代码实现}
    \end{minipage}
\end{figure}

\begin{lstlisting}[language=Python]
# 加载预训练的 RNA-FM 模型
# model_location: 预训练权重文件的本地路径
self.backbone, self.alphabet = fm.pretrained.rna_fm_t12(
    model_location=rna_fm_model_path
)

# 提取第 12 层的隐藏状态作为特征,repr_layers=[12] 指定提取层数
with torch.no_grad():
    results = self.backbone(tokens, repr_layers=[12])

# 获取特征张量，维度为 [批次大小, 序列长度, 640]
embedding = results["representations"][12]
\end{lstlisting}

\subsection*{3.2 维度变换与成对特征构建}

\hspace{2em}本实验的核心挑战在于：如何将一维的序列特征转化为二维的结构特征？因为 RNA 的二级结构本质上是一个 $L \times L$ 的碱基配对矩阵，而 RNA-FM 输出的是 $L \times D$ 的序列向量。为了解决这一维度跨越问题，本实验设计了外积拼接模块。

该模块利用广播机制，显式地构建了序列中任意两个位置 $i$ 和 $j$ 之间的相互关系。具体实现步骤如下：

\begin{itemize}
    \item \textbf{行特征广播}：
    
    \hspace{2em}将输入特征矩阵 $X$ 在“列”维度上进行复制扩展。对于输出张量的每一个位置 $(i, j)$，填充第 $i$ 个碱基的特征向量。这在几何上表现为将 $L \times D$ 的特征板水平拉伸，形成具有“水平条纹”特征的立方体张量。
    \item \textbf{列特征广播}：
    
    \hspace{2em}将输入特征矩阵 $X$ 在“行”维度上进行复制扩展。对于输出张量的每一个位置 $(i, j)$，填充第 $j$ 个碱基的特征向量。这在几何上表现为将 $L \times D$ 的特征板垂直拉伸，形成具有“垂直条纹”特征的立方体张量。
    \item \textbf{特征融合}：
    
    \hspace{2em}将上述两个张量在通道维度进行拼接。最终生成的特征体 $Z$ 中的每一个点 $(i, j)$ 都同时包含了第 $i$ 个碱基和第 $j$ 个碱基的特征信息，特征维度翻了一倍。
\end{itemize}

\begin{lstlisting}[language=Python]
# 输入 x 形状: [B, L, D]

# 1. 扩展维度并复制，构建"行"特征
# x.unsqueeze(2) -> [B, L, 1, D]
# .expand -> [B, L, L, D]
x_row = x.unsqueeze(2).expand(-1, -1, L, -1)

# 2. 扩展维度并复制，构建"列"特征
# x.unsqueeze(1) -> [B, 1, L, D]
# .expand -> [B, L, L, D]
x_col = x.unsqueeze(1).expand(-1, L, -1, -1)

# 3. 在最后一个维度(特征维)进行拼接
# 最终形状: [B, L, L, 2D]
out = torch.cat((x_row, x_col), dim=-1)
\end{lstlisting}

\begin{figure}[H]
    \centering
    \begin{minipage}{\linewidth}
        \centering
        % 请确保图片路径正确
        \includegraphics[width=\textwidth]{代码实现/concat.png} 
        \caption{Outer Concatenation 维度变换代码实现}
    \end{minipage}
\end{figure}

\newpage

\subsection*{3.3 结构预测与解码}

\hspace{2em}本模块是模型的核心决策层，负责从高维特征中解析出稀疏的配对模式，并输出最终的概率矩阵。整个解码过程分为降维预处理、深度特征提取和概率输出三个阶段。

\begin{enumerate}
  \item \textbf{降维预处理}:
  
  \hspace{2em}为了降低计算复杂度，模型首先引入了一个全连接层作为瓶颈层。由于上一阶段生成的特征体通道数高达 1280 维（2D），直接进行卷积运算开销过大。该层将特征通道数大幅压缩至 128 维（记为H，为放缩后的维度），在保留关键结构信息的同时显著减少了参数量。随后，通过维度置换操作，将数据调整为适配二维卷积的格式（[B, L, L, H] -> [B, H, L, L]）。
  \item \textbf{深度特征提取}:
  
  \hspace{2em}模型采用了包含 16 个模块的深度残差网络（Resnet）。与传统的卷积网络不同，本项目采用了“瓶颈式”的三层卷积结构，包含$ 1 \time 1 \rightarrow 3 \time 3 \rightarrow 1 \time 1$卷积。层间均使用Relu函数进行激活。此外，残差连接机制贯穿每个模块，确保了梯度在深层网络中的有效传播，使模型能够捕捉长距离的碱基依赖关系。
  \item \textbf{概率输出}:
  
  \hspace{2em}最后，经过深层提取的特征被送入输出层。该层通过一个卷积核大小为 3 的二维卷积层，将 128 个特征通道压缩为 1 个通道，得到最终的配对得分矩阵。在推理阶段，该矩阵的数值通过激活函数映射到 0 到 1 区间，表示任意两个碱基形成配对的置信度。
\end{enumerate}

\begin{figure}[H]
    \centering
    \begin{minipage}{0.85\linewidth}
        \centering
        % 请在 models/my_model.py 中截取 ClassifierHead 类的代码
        \includegraphics[width=\textwidth]{代码实现/ClassifierHead_init.png}
        \caption{ClassifierHead的部分代码}
    \end{minipage}
\end{figure}

\begin{lstlisting}[language=Python]
# 1. 降维层：将 1280 维的高维特征压缩为 128 维
self.linear_in = nn.Linear(2 * in_dim, hidden_dim)
# 2. 特征提取器：堆叠 16 层残差网络
self.resnet = ResNet2D(hidden_dim, num_blocks, kernel_size)
# 3. 输出层：将特征通道压缩为 1，输出配对得分
self.conv_out = nn.Conv2d(hidden_dim, 1, kernel_size=kernel_size, padding='same')

# 前向传播流程
x = self.outer_concat(x)      # 维度变换
x = self.linear_in(x)         # 线性降维
x = x.permute(0, 3, 1, 2)     # 调整维度顺序以适配卷积
x = self.resnet(x)            # 深度残差提取
x = self.conv_out(x)          # 输出最终预测
\end{lstlisting}



\subsection*{3.4 运行结果截图}

\hspace{2em}由于训练时间很长，平均每个Epoch需要约2小时。因此我选择设计一段可以从前一个Epoch开始训练的代码，也因此我的输出结果并不能展示整个训练过程。但是我将运行结果写入了日志，下面将以日志截图代替训练的输出结果截图。忽略其中的Acc，因为我一开始计算了Acc，之后为了节省时间没有计算Acc指标，将其令为0。最终效果可见上传的train.log文件，在这里仅展示训练了1个epoch的模型
接着训练的命令行截图。

\begin{figure}[H]
    \centering
    \includegraphics[width=0.65\textwidth]{logs/image.png}
    \caption{训练1个epoch后接着训练的截图}
    \label{fig:log}
\end{figure}

\newpage

% =========================================================
% 四、实验结果 (严格复刻要求，且表格使用三线表)
% =========================================================
\section*{四、实验结果}

\subsection*{4.1 模型参数及结构描述}

\hspace{2em}本实验提出的 RNA 二级结构预测模型采用了“预训练语言模型 + 深度残差网络”的混合架构。整体结构如图 \ref{fig:model_arch} 所示，主要包含以下三个核心处理阶段：

\begin{enumerate}
    \item \textbf{序列特征嵌入 (Embedding Stage)}：
    
    \hspace{2em}利用预训练的 \textbf{RNA-FM} 模型作为特征提取器。对于长度为 $L$ 的输入 RNA 序列，RNA-FM 输出包含RNA序列信息的嵌入向量，其维度为 $L \times D$。为节省计算资源并防止过拟合，本实验冻结了 RNA-FM 的骨干参数。

    \item \textbf{成对特征构建 (Pairwise Representation Construction)}：
    
    \hspace{2em}为了将一维序列信息转化为矩阵，模型引入了外部拼接机制。如图\ref{fig:model_arch}所示，模型分别在“行”维度和“列”维度对特征进行广播（Broadcasting）。随后将两者在通道维度进行拼接，生成形状为 $[L, L, 2D]$ 的二维特征张量。这一步使得每个空间位置 $(i, j)$ 都包含了第 $i$ 个碱基和第 $j$ 个碱基的组合特征。

    \item \textbf{深度特征提取与预测 (Deep Feature Extraction)}：
    
    \hspace{2em}生成的二维特征图被送入包含 16 个残差块（Residual Blocks）的 \textbf{ResNet} 中。利用卷积层和残差连接，模型能够有效捕捉长距离的碱基依赖关系。最终，通过一个 $1 \times 1$ 的卷积层将通道压缩，并经过 Sigmoid 激活函数输出大小为 $L \times L$ 的概率矩阵，表示碱基配对的置信度。
\end{enumerate}

\begin{figure}[H]
  \centering
  % 这里使用了你提供的图片路径，请确保文件名无误
  \includegraphics[width=0.8\textwidth]{代码实现/原理图.png}
  \caption{基于 RNA-FM 与 ResNet 的 RNA 二级结构预测模型架构示意图。}
  \label{fig:model_arch}
\end{figure}

\subsection*{4.2 训练参数}

\hspace{2em}本实验的具体训练超参数设置如表 \ref{tab:hyperparams} 所示。为了解决正负样本（配对与非配对）极度不平衡的问题，损失函数中引入了正样本权重（pos\_weight）。

\begin{table}[H]
    \centering
    \caption{实验模型参数与训练超参数设置}
    \label{tab:hyperparams}
    \begin{tabular}{lc}
        \toprule
        \textbf{参数名称 (Parameter)} & \textbf{数值 (Value)} \\
        \midrule
        预训练模型 (Pre-trained Model) & RNA-FM (冻结参数) \\
        放缩后的维度 (Hidden, $H$) & 128 \\
        ResNet 层数 (Blocks) & 16 \\
        最大序列长度 (Max Length, $L$) & 1022 \\
        优化器 (Optimizer) & AdamW \\
        学习率 (Learning Rate) & $1 \times 10^{-4}$ \\
        批大小 (Batch Size) & 2 \\
        梯度累积步数 (Accumulation Steps) & 4 \\
        损失函数 (Loss Function) & 二元交叉熵损失  \\
        正样本权重(pos\_weight) & 10 \\
        训练轮数 (Epochs) & 10 \\
        \bottomrule
    \end{tabular}
\end{table}

\subsection*{4.3 评估指标}
% (在此处分析结果)

\subsubsection*{4.3.1 训练与验证：RNAStrAlign}
\hspace{2em}为了评估模型在训练过程中的收敛情况及泛化能力，我们绘制了训练集与验证集的损失变化曲线（图 \ref{fig:loss}）以及验证集上的关键性能指标变化曲线（图 \ref{fig:metrics}）。

\begin{figure}[H]
    \centering
    \begin{minipage}[b]{0.48\textwidth}
        \centering
        % 请确保图片路径正确，建议尽量使用英文路径以避免编译错误
        \includegraphics[width=\textwidth]{评估/loss_curve.png}
        \caption{训练与验证损失收敛曲线}
        \label{fig:loss}
    \end{minipage}
    \hfill
    \begin{minipage}[b]{0.48\textwidth}
        \centering
        \includegraphics[width=\textwidth]{评估/metrics_curve.png}
        \caption{验证集中 P, R 与 F1 变化曲线}
        \label{fig:metrics}
    \end{minipage}
\end{figure}

结果分析：
\begin{itemize}
    \item \textbf{损失收敛性}：
    
    \hspace{2em}如图 \ref{fig:loss} 所示，Epoch≤8时，训练损失与验证损失均呈现下降趋势。但是之后验证损失呈现上升趋势，说明模型在训练过程中出现了过拟合现象，不过并不严重。
    \item \textbf{指标变化}：
    
    \hspace{2em}如图 \ref{fig:metrics} 所示，F1 分数随训练过程稳步提升。模型在保持较高召回率的同时，精确率逐步改善，表明模型区分真假阳性配对的能力在不断增强。
\end{itemize}

\subsubsection*{4.3.2 基准测试：ArchiveII}
\hspace{2em}在模型训练完成后，我们在独立的测试集（Benchmark Dataset）上进行了最终性能评估。该测试集的数据分布与训练集隔离，能够客观反映模型的泛化性能。最终的评估结果如表 \ref{tab:test_results} 所示。

\begin{table}[H]
    \centering
    \caption{测试集最终性能评估结果}
    \label{tab:test_results}
    \vspace{0.5em}
    \begin{tabular}{lccccc}
        \toprule
        \textbf{数据集} & \textbf{准确率} & \textbf{精确率} & \textbf{召回率} & \textbf{F1} & \textbf{AUC} \\
        \midrule
        ArchiveII & 0.9978 & 0.7148 & 0.7432 & 0.7287 & 0.9632 \\
        \bottomrule
    \end{tabular}
\end{table}

结果分析：
\begin{enumerate}
    \item \textbf{极高的准确率与稀疏性}：
    
    \hspace{2em}实验结果显示，模型的准确率高达 \textbf{99.78\%}。虽然很高，但是RNA 二级结构接触图具有极高的稀疏性，因为绝大多数碱基不发生配对，即负样本占绝大多数。因此，准确率主要反映了模型对负样本（TN）的正确排除能力，而非单一的核心性能指标。
    \item \textbf{综合性能}：
    
    \hspace{2em}F1 分数为 \textbf{0.7287}，该指标表明模型在正样本预测上取得了良好的平衡，既能以较高的比例覆盖真实配对，又能将误报率控制在合理范围内。
    \item \textbf{AUC}：
    
    \hspace{2em}AUC 达到了 \textbf{0.9632}，这表明模型输出的预测概率具有极高的可信度。在不同的分类阈值下，模型能够极其有效地将真实配对与非配对区分开来，证明了模型架构在特征提取方面的有效性。
\end{enumerate}

\subsection*{4.4 ROC曲线}

\begin{figure}[H]
    \centering
    \includegraphics[width=0.9\textwidth]{评估/test_roc_curve.png}
    \caption{模型在测试集上的ROC曲线}
    \label{fig:roc_curve}
\end{figure}

\newpage

% =========================================================
% 五、系统实现与交互
% =========================================================
\section*{五、前端界面展示}
% (在此处展示UI截图等)
\subsection*{5.1 首页}

首页页面过长，为了报告的美观，仅展示部分，具体见README.md

\begin{figure}[H]
    \centering
    \includegraphics[width=\textwidth]{UI/首页.png}
    \caption{首页预览图}
    \label{fig:home}
\end{figure}

同时，可以和右侧的阿里云AI对话，该AI通过检索增强技术（RAG），获得了我所有的代码文件，可以回答大多数问题：

\begin{figure}[H]
    \centering
    % 左边的图片 (占用 48% 宽度)
    \begin{minipage}[b]{0.48\textwidth}
        \centering
        % 请确保文件名和路径正确，LaTeX 中建议使用斜杠 /
        \includegraphics[width=\textwidth]{UI/bot_1.png}
        \caption{AI 对话界面展示 (1)}
        \label{fig:bot1}
    \end{minipage}
    \hfill % 这里的 \hfill 是关键，它负责把两张图撑开到两边
    % 右边的图片 (占用 48% 宽度)
    \begin{minipage}[b]{0.48\textwidth}
        \centering
        \includegraphics[width=\textwidth]{UI/bot_2.png}
        \caption{AI 对话界面展示 (2)}
        \label{fig:bot2}
    \end{minipage}
\end{figure}

\subsection*{5.2 单条序列}

在未填入序列的情况下：

\begin{figure}[H]
    \centering
    \includegraphics[width=\textwidth]{UI/单条_未填入.png}
    \caption{单条序列-未填入}
\end{figure}

填入序列后的预测结果：

\begin{figure}[H]
    \centering
    \includegraphics[width=\textwidth]{UI/单条_已填入.png}
    \caption{单条序列-已填入}
\end{figure}

若目前没有可用的fasta序列，可以选择训练集里RNA家族中的某一个示例，如图所示：

\begin{figure}[H]
    \centering
    % 左边的图片 (占用 48% 宽度)
    \begin{minipage}[b]{0.48\textwidth}
        \centering
        % 请确保文件名和路径正确，LaTeX 中建议使用斜杠 /
        \includegraphics[width=\textwidth]{UI/单条_RNA.png}
        \caption{单条序列-选择RNA家族}
    \end{minipage}
    \hfill % 这里的 \hfill 是关键，它负责把两张图撑开到两边
    % 右边的图片 (占用 48% 宽度)
    \begin{minipage}[b]{0.48\textwidth}
        \centering
        \includegraphics[width=\textwidth]{UI/单条_fasta.png}
        \caption{单条序列-选择fasta}
    \end{minipage}
\end{figure}

\subsection*{5.3 批量预测}

在批量预测页面，可以上传多个fasta文件进行预测：

\begin{figure}[H]
    \centering
    \includegraphics[width=\textwidth]{UI/批量_未填入.png}
    \caption{批量预测-未填入}
\end{figure}

上传文件后的界面：

\begin{figure}[H]
    \centering
    \includegraphics[width=\textwidth]{UI/批量_已填入.png}
    \caption{批量预测-已填入}
\end{figure}

\subsection*{5.4 验证对比}

在该页面，可以选择带有真实结构的bpseq文件进行对比验证：

\begin{figure}[H]
    \centering
    \includegraphics[width=\textwidth]{UI/对比_未填入.png}
    \caption{验证对比-未填入}
\end{figure}

同样，可以选择不同RNA家族的不同示例进行对比，包含训练集和测试集：

\begin{figure}[H]
    \centering
    % 左边的图片 (占用 48% 宽度)
    \begin{minipage}[b]{0.48\textwidth}
        \centering
        % 请确保文件名和路径正确，LaTeX 中建议使用斜杠 /
        \includegraphics[width=\textwidth]{UI/对比_RNA.png}
        \caption{验证对比-选择RNA家族}
    \end{minipage}
    \hfill % 这里的 \hfill 是关键，它负责把两张图撑开到两边
    % 右边的图片 (占用 48% 宽度)
    \begin{minipage}[b]{0.48\textwidth}
        \centering
        \includegraphics[width=\textwidth]{UI/对比_bpseq.png}
        \caption{验证对比-选择bpseq}
    \end{minipage}
\end{figure}

\newpage
% =========================================================
% 六、总结与心得
% =========================================================
\section*{六、经验及总结}

\subsection*{6.1 分点总结}

除了第一条外，我将简单列举我遇到的问题与解决方法：

\begin{enumerate}
    \item 最严重、代价最大
    \begin{itemize}
        \item 所遇问题：我在第一次训练时，当时还没出大作业模板，由于我的训练意识比较薄弱，没有考虑到使用固定种子去划分验证集和训练集，导致我最后发现我无法找到验证集去计算指标。
        \item 解决方法：我重新设计了增加计算准确率的函数，并且稍微改了下训练代码ssp.py，然后使用了全局种子42，最终得以解决该问题。
        \item 自我总结：虽然我最后重新解决了该问题，计算出了每一个Epoch的四个指标，但是浪费了我的阿里云AI工作站的算力。最后我训练了10个Epoch，消耗了我的大量算力，本来我还打算留不少算力用于训练《医工交叉融合的人工智能》里的大作业，现在还得找别人的账号或者充钱了。
    \end{itemize}
    \item 追求完美
    \begin{itemize}
        \item 所遇问题：在第二次训练的时候，由于为了追求获得更多的指标，我决定修改计算指标的metrics函数，在每次Epoch跑完后都对验证集同时计算准确率、精确率、召回率和F1分数。但是，由于计算准确率太复杂，导致在验证集的评估上会花费大量时间，浪费了我的算力。
        \item 解决方法：考虑到目前有限的算力，我不得不放弃计算过程中对验证集的准确率。
        \item 自我总结：有的时候，追求完美并不是一件好事，尤其是在算力有限的情况下，更何况在样本种类失衡的状况下，准确率显然不如另外三个指标重要。我以后在训练模型时，会优先考虑算力的消耗，避免不必要的计算。
    \end{itemize}
    \item 数据加载
    \begin{itemize}
        \item 所遇问题：我在加载数据时，会报错文件不存在，最后发现该函数只能进入一个文件夹，不能进入子文件夹并检索数据。
        \item 解决方法：我修改了加载数据的函数，使用 os.walk 替代 glob，os.walk 会自动进入所有子文件夹，解决 RNAStrAlign 的嵌套问题。
        \item 自我总结：以后在加载数据时，我会检查数据文件夹的结构，如果存在子文件夹，我会优先考虑使用 os.walk 进行递归检索。
    \end{itemize}
    \item 搬运代码
    \begin{itemize}
        \item 所遇问题：由于我本地的电脑显存仅有6GB，而训练该模型需要很高的算力，因此我决定租一个阿里云的AI工作站，在上面进行训练。可是后面我发现，重新配置环境极其麻烦，因为其网络有限，很多包无法下载，而只要使用就会消耗算力时间。同时，我的数据过大，搬运到工作站上也很麻烦。
        \item 解决方法：数据方面，我向客服发出工单，客服告诉我如何高速传输数据；同时，我将环境压缩成压缩包，在工作站里访问本地的文件夹，将本地的环境压缩包在工作站里解压，这样就可以不依赖于网络就快速配置好环境了。
        \item 自我总结：幸好AI工作站可以直接访问我本地的文件夹，否则我可能需要花费非常非常多的时间去配置环境，以至于放弃使用AI工作站。
    \end{itemize}    
    \item 显存爆炸
    \begin{itemize}
        \item 所遇问题：在训练过程中，我使用的是阿里云48G显存的GPU，但是在训练的时候，要么就是显存不够，报错CUDA out of memory，要么就是显存利用率少，上监控一看我的显存利用率才30\%，浪费我租工作站的钱。
        \item 解决方法：经过一晚上的调试，我终于找到一个合适的超参数，使得显存不会爆掉，同时显存利用率也能达到85\%以上。我将\texttt{batch\_size}调小为2，同时将模型的残差网络层数调小为16层，最终达到了一个折中的效果。
        \item 自我总结：在训练过程中，我使用GPU训练时，将\texttt{batch\_size}调小，使得显存不会爆掉。
    \end{itemize}
    \item RAG实现
    \begin{itemize}
        \item 所遇问题：在搭建 WebUI 的“智能问答助手”模块时，我遭遇了 Gradio 版本更新带来的巨大阵痛。原本参考教程编写的 \texttt{rag\_chat} 函数在运行时疯狂报错。查阅文档后才发现，新版 Gradio 的 Chatbot 组件已经弃用了元组格式，强制要求使用字典列表格式。
        \item 解决方法：我不得不重写了整个问答逻辑，适配了最新的 \texttt{messages} 格式。此外，我编写了一个后端函数，在启动时自动递归扫描我指定的文件并注入到系统提示词中，实现了RAG。
        \item 自我总结：开源社区更新太快了，教程往往是滞后的。在做工程实现时，不能死板照抄代码，必须具备阅读最新官方文档的能力。
    \end{itemize}
    \item 追求美观
    \begin{itemize}
        \item 所遇问题：为了追求美观，我用Nano banana Pro绘制了背景图，很有RNA结构的生物科技感，并希望按钮颜色能吸取背景中的主色调，但在实际操作中发现，我实在找不到如何更改可视化界面的背景图，且按钮的默认色调极难通过简单参数覆盖。
        \item 解决方法：我利用取色器提取了背景图中的关键色值，不再依赖默认主题，而是通过 \texttt{CSS} 注入的方式，强制重写了 \texttt{.gradio-container} 的背景属性和 \texttt{button.primary} 的颜色样式。经过多次尝试，我终于更改了按钮的颜色，更换了整体的主题。
        \item 自我总结：虽然为了调这几个颜色花费了很多时间，甚至背景图最终还是没有整出来，但我在这其中学到的不仅是如何更换主题颜色，还学到了如何上网搜索教程以满足自己的美观需求。
    \end{itemize}
    \item 急于求成
    \begin{itemize}
        \item 所遇问题：我的模型在验证集上跑的时候，用时非常长，足足40分钟。于是我尝试修改其\texttt{batch\_size}，从2增加到16。结果吃完饭回来，发现显存已经爆了，我白白浪费了一个半小时的算力。之后，我将\texttt{batch\_size}改成4，发现显存占用从85\%降低到了约60\%，反而降低了速度。
        \item 解决方法：我将\texttt{batch\_size}改回来了。
        \item 自我总结：我太激进了，同时太冲动了。我忘记了我的保存模型代码是在验证评估之后，因此一旦显存爆了，那我的模型就没了。同时，我也只是想当然的改了\texttt{batch\_size}，结果反而又更加耗时了。这一连串操作一共让我白白浪费了快2个小时，都够我训练一个Epoch+验证了。我的心急反而让我付出了更多的代价，实在糊涂！不过说实话，我到现在都不知道为什么我的\texttt{batch\_size}变大反而显存占用更少了，导致速度更慢了。等放假我要好好查明原因，避免下次再犯错。
    \end{itemize}
    \item 自动运行
    \begin{itemize}
        \item 所遇问题：我刚开始尝试运行的时候，总是在不知不觉间爆显存，浪费了算力。
        \item 解决方法：我写了一个\texttt{run\_with\_restart.py}，可以在显存爆了后自动重新运行。
        \item 自我总结：虽然感觉没能解决根本问题，但至少获得了一些经验，而且之后也没有爆过显存。
    \end{itemize}
\end{enumerate}

\subsection*{6.2 个人总结}

\hspace{2em}在该大作业中，我实现了第一次独立开发一个AI小项目，并且与我的专业高度相关。通过该实验，我亲自完成了从思路构建、代码实现、模型训练到设计UI的完整流程，收获颇丰。为了锦上添花，我决定去获取阿里云的免费限额API，通过RAG的技术接入到我的UI界面中，这样在展示的时候别人也可以通过该AI助手，快速简单了解我的项目。

由于至始至终都是一个人，我不得不在冬学期初就开始筹划如何开展。刚开始时，是我认识的师兄为我提供了灵感，他之前做过类似的项目，并且给我推荐了RNA-FM这个预训练模型，让我设计新的概率预测算法。随后，我花了一周时间去熟悉该模型的使用方法，并且阅读了相关论文，了解其原理。接下来，我开始设计我的模型结构，并且编写代码实现。在这个过程中，我遇到了很多困难，比如如何进行维度变换，如何设计残差网络等。通过查阅资料和反复调试，我最终解决了这些问题。

同时，也很感谢朱老师。因为一开始我在学期初就和别人组好队了，但是朱老师很看好我，认为我可以自己一个人做，因此我才下定决心，离开队伍后自行开发。同时，朱老师还为我提供了一些科研经费，让我不那么算力焦虑。不过，至少目前我还可以白嫖阿里云的算力，也很感谢阿里云的算力支持，让我有能力训练模型。我计算了一下，我一共耗时大概3200核时，折合人民币大约要100元。如果没有阿里云大方送我们算力的话，我是肯定舍不得这么多钱投入到我的作业当中的。

\end{document}